\documentclass[fleqn]{report}

% Neðangreindar þrjá línur eru nauðsynlegar fyrir íslensku, t.d. til að geta skrifað Þ og Æ.
\usepackage[icelandic]{babel}
\usepackage{tikz}
\usepackage[T1]{fontenc}
\usepackage{amssymb}
\usepackage[utf8]{inputenc}
\usepackage[fleqn]{amsmath}
\usepackage{amssymb}
\usepackage{amsfonts}
\usepackage{mathtools}
\usepackage{float}
\usepackage{caption}
\usepackage{tcolorbox}
\usepackage{setspace}
\usepackage{sectsty}
\usepackage{hyperref}

\chapterfont{\centering\scshape\mdseries}
\sectionfont{\centering\scshape\mdseries}
\subsectionfont{\centering\scshape\mdseries}

\onehalfspacing

\usepackage{newtxtext, newtxmath}

\usepackage{caption}
\DeclareCaptionLabelSeparator{pkt}{. }
\captionsetup{
    labelfont = {bf, it},
    labelsep  = pkt,
    font = small
}

\usepackage{booktabs}

\usepackage{parskip}
\usepackage{enumerate}
\graphicspath{{./img/}}
\usepackage[a4paper,top=2cm,bottom=2cm,left=3cm,right=3cm,marginparwidth=1.75cm]{geometry}

% pakki fyrir heimildavinnu (þarf yfirleitt ekki fyrir heimadæmi)
\usepackage[numbers]{natbib}

% graphicx er pakki sem er m.a. notaður fyrir myndir
\usepackage{graphicx}

% Booktabs er pakki sem gerir töflur snyrtilegri
\usepackage{booktabs}

% Pakki sem gerir okkur kleyft að nota align umhverfið
\usepackage{amsmath}
\usepackage{enumerate}
\usepackage{listings}
\usepackage{color}
\usepackage{bm}
\usepackage{tikz}

\newcommand{\ms}{\hspace{\minuslength}}

\definecolor{dkgreen}{rgb}{0,0.6,0}
\definecolor{gray}{rgb}{0.5,0.5,0.5}
\definecolor{mauve}{rgb}{0.58,0,0.82}

\definecolor{codegreen}{rgb}{0,0.6,0}
\definecolor{codegray}{rgb}{0.5,0.5,0.5}
\definecolor{codepurple}{rgb}{0.58,0,0.82}
\definecolor{backcolour}{rgb}{0.95,0.95,0.92}

\usepackage{wasysym}
\font\wasy=wasy10 at 10pt
\def\thorn{{\wasy\char105}}
\def\eth{{\wasy\char107}}

\def\lstlistingname{Listun}
\lstset{
  language = Python, 
  frame = single, 
  escapeinside={(*@}{@*)},
  aboveskip=2mm,
  belowskip=1mm,
  showstringspaces=false,
  columns=flexible,
  basicstyle={\small\ttfamily},
  numbers=left,
  numberstyle=\tiny\color{gray},
  keywordstyle=\color{blue},
  commentstyle=\color{dkgreen},
  stringstyle=\color{mauve},
  breaklines=true,
  breakatwhitespace=true,
  tabsize=3,
  captionpos=b,
  literate=%
    {é}{{\'{e}}}1
    {è}{{\`{e}}}1
    {ê}{{\^{e}}}1
    {ë}{{\¨{e}}}1
    {É}{{\'{E}}}1
    {Ê}{{\^{E}}}1
    {û}{{\^{u}}}1
    {ù}{{\`{u}}}1
    {ú}{{\'{u}}}1
    {â}{{\^{a}}}1
    {à}{{\`{a}}}1
    {á}{{\'{a} }}1
    {ã}{{\~{a}}}1
    {Á}{{\'{A}}}1
    {Â}{{\^{A}}}1
    {Ã}{{\~{A}}}1
    {ç}{{\c{c}}}1
    {Ç}{{\c{C}}}1
    {õ}{{\~{o}}}1
    {ó}{{\'{o}}}1
    {ô}{{\^{o}}}1
    {Õ}{{\~{O}}}1
    {Ó}{{\'{O}}}1
    {Ô}{{\^{O}}}1
    {î}{{\^{i}}}1
    {Î}{{\^{I}}}1
    {í}{{\'{i}}}1
    {Í}{{\~{Í}}}1
    {ð}{{\eth}}1
}

\newlength{\minuslength}
\settowidth{\minuslength}{$-$}
\newcommand{\R}{\mathbb{R}}
\newcommand{\T}{\mathcal{T}}
\newcommand{\N}{\mathbb{N}}
\newcommand{\titill}{Réttarvísir}

\setcounter{secnumdepth}{0}

% Titill
\title{
\includegraphics[width = 0.45\textwidth]{Tákn blátt.pdf}
\vspace{2cm}

\textsc{TÖL025M Inngangur að máltækni} \\ 
\Huge \titill}
\author{\Large Jóhannes Reykdal Einarsson\\\Large Sölvi Santos}

\date{\Large\today \vfill \large Lokaverkefni í TÖL025M \\ Verk- og náttúruvísindasvið}

\begin{document}
\maketitle

\tableofcontents
\thispagestyle{empty}
\newpage

\setcounter{page}{1}

\chapter*{\titill}
\addcontentsline{toc}{chapter}{\titill}
\section{Inngangur}

\section{Bakgrunnur}

\section{Umræður}

\section{Næstu skref}

\section{Hvað fór úrskeiðis}

\section{Samstarf}


\chapter*{Skjölun}
\addcontentsline{toc}{chapter}{Skjölun}
\section{Bakendi}


\newpage

\appendix

\chapter*{Viðauki}
\addcontentsline{toc}{chapter}{Viðauki}
\section{Skilgreiningar}
\fbox{%
  \parbox{\linewidth}{%
    \textbf{Skilgreining 1}\\
    Látum $f_{t,d}$ tákna hversu oft hugtak $t$ kemur fyrir í skjali $d$. Tíðni hugtaks (e. term frequency) í skjali $d$ er skilrgeint sem hlutfallsleg tíðni hugtaks í skjali,\begin{align*}
        \text{tf}(t,d) = \frac{f_{t,d}}{\sum_{t'\in d}f_{t',d}}.
    \end{align*}
  }%
}

\fbox{%
  \parbox{\linewidth}{%
    \textbf{Skilgreining 2}\\
    Hallo
  }%
}

\newpage

\addcontentsline{toc}{chapter}{Heimildaskrá}
\begin{thebibliography}{99}
    \bibitem{rna-seq}
    Kukurba, K. R., \& Montgomery, S. B. (2015). \textit{RNA sequencing and analysis}. Cold Spring Harbor Protocols.

    \bibitem{transcriptome}
Caudai, C., Galizia, A., Geraci, F., Le Pera, L., Morea, V., Salerno, E., Via, A., \& Colombo, T. (2021). \textit{AI applications in functional genomics}. Computational and Structural Biotechnology Journal, 19, 5762–5790.

\bibitem{fastq}
Cock, P. J. A., Fields, C. J., Goto, N., Heuer, M. L., \& Rice, P. M. (2010).
\textit{The Sanger FASTQ file format for sequences with quality scores, and the Solexa/Illumina FASTQ variants}. Nucleic Acids Research 38(6):1767-71.    

\bibitem{wass}
Wasserman, L. (2004). \textit{All of statistics: A concise course in statistical inference}. Springer.

\bibitem{lag}
Lagrange multiplier. (n.d.). In Wikipedia. Retrieved March 29, 2026, from \href{https://en.wikipedia.org/wiki/Lagrange_multiplier}{https://en.wikipedia.org/wiki/Lagrange\_multiplier}

\bibitem{kallisto}
Bray, N. L., Pimentel, H., Melsted, P., \& Pachter, L. (2016). \textit{Near-optimal probabilistic RNA-seq quantification}. Nature Biotechnology, 34(5), 525–527.

\bibitem{equiv}
? \textit{Stopporð}. Retrieved March 30, 2026, from \url{https://github.com/atlijas/icelandic-stop-words/blob/master/all_stop_words.txt}


\end{thebibliography}

\end{document}
